% Generated draft: requires mathematical and source review.
\documentclass[11pt]{article}
\usepackage[margin=25mm]{geometry}
\usepackage{fontspec}
\setmainfont{Latin Modern Roman}
\usepackage{amsmath,amssymb,mathtools,unicode-math}
\setmathfont{Latin Modern Math}
\usepackage{enumitem,needspace,xurl,hyperref}
\setlength{\parindent}{0pt}
\setlength{\parskip}{5pt}
\setlength{\emergencystretch}{3em}
\newcommand{\N}{\mathbb{N}}
\newcommand{\Z}{\mathbb{Z}}
\newcommand{\Q}{\mathbb{Q}}
\newcommand{\R}{\mathbb{R}}
\newcommand{\C}{\mathbb{C}}
\newcommand{\F}{\mathbb{F}}
\newcommand{\A}{\mathbb{A}}
\newcommand{\PP}{\mathbb{P}}
\newcommand{\E}{\mathbb{E}}
\newcommand{\eps}{\varepsilon}
\newcommand{\dd}{\,\mathrm{d}}
\begin{document}
\section*{UnsolvedMath ID: 30003793}
\textbf{Source problem code:} catalogue entry
\subsection{Definitions and mathematical statement}
Let \(R\) be a discrete valuation ring, let \(F=\operatorname{Frac}(R)\), and let \(k\) be its residue field. The source requires that \(2\) be invertible, but its formulations do not consistently locate this condition: \(2\in R^\times\) and \(2\in F^\times\) give different hypotheses. The intended choice remains unresolved below.

An \(R\)-order \(A\) is a unital \(R\)-algebra that is finitely generated and torsion-free as an \(R\)-module and spans its finite-dimensional generic algebra \(A_F=A\otimes_R F\). It is hereditary if every left ideal is projective as a left \(A\)-module. In the usual \(R\)-linear interpretation, an involution is an \(R\)-linear map \(\tau:A\to A\) satisfying
\[
\tau(ab)=\tau(b)\tau(a),\qquad \tau(1)=1,\qquad \tau^2=\operatorname{id}_A.
\]
An isomorphism of orders with involution is an \(R\)-algebra isomorphism \(f:A\to A'\) such that \(f\circ\tau=\tau'\circ f\). For an \(R\)-algebra \(T\), write \((A_T,\tau_T)=(A\otimes_R T,\tau\otimes\operatorname{id}_T)\).

Let \(J(A)\) denote the Jacobson radical. An involution preserves this radical and therefore induces an involution \(\overline\tau\) on \(\overline A=A/J(A)\). Here anisotropy means
\[
\forall x\in\overline A,\qquad \overline\tau(x)x=0\ \Longrightarrow\ x=0.
\]
Use the analogous notation and definition for \((A',\tau')\).

\textbf{Principal descent question.} For every pair of hereditary \(R\)-orders with involution \((A,\tau)\) and \((A',\tau')\), suppose that
\[
(A_F,\tau_F)\cong(A'_F,\tau'_F)
\]
and that there exists a finite \'{e}tale \(R\)-algebra \(B\) for which
\[
(A_B,\tau_B)\cong(A'_B,\tau'_B).
\]
If both \(\overline\tau\) and \(\overline\tau'\) are anisotropic, must
\[
(A,\tau)\cong(A',\tau')
\]
over \(R\)? All displayed isomorphisms are isomorphisms of algebras with involution. The source does not explicitly require \(B\) to be nonzero or faithfully flat; this qualification must be settled to specify the intended \'{e}tale descent hypothesis. Over a discrete valuation ring, a nonzero finite \'{e}tale algebra is faithfully flat.

\textbf{Further questions, incompletely specified by the source.} Does the same descent conclusion hold when the residual-anisotropy hypothesis is replaced by the condition that every closed fiber of the automorphism group scheme is \emph{almost anisotropic}? The natural automorphism functor of a pair is
\[
T\longmapsto\operatorname{Aut}_{T\text{-alg},\tau_T}(A_T),
\]
but the source does not specify whether the condition concerns this functor for \(A\), the corresponding functor for \(A'\), or both. Nor does it define ``almost anisotropic.'' For a group scheme over \(\operatorname{Spec}R\), there is only one closed base point, with fiber over \(k\); whether a geometric closed fiber is intended is also unspecified.

Finally, the source asks whether every central simple \(F\)-algebra with involution \((C,\sigma)\) admits a hereditary \(R\)-order with involution whose generic fiber is isomorphic to \((C,\sigma)\) and which is ``such an order.'' It does not specify whether this phrase requires residual anisotropy, almost anisotropy of an automorphism-group fiber, or another preceding property. The permitted action of involutions on the center is likewise unspecified. Consequently this existence question, and the proposed group-scheme generalization, cannot yet be given a unique precise formulation.
\subsection{Short English statement}
The principal question asks whether two hereditary orders with involution that are isomorphic over the fraction field and after a finite \'{e}tale extension must already be isomorphic over the discrete valuation ring when their radical quotients have anisotropic involutions. The source also asks for an almost-anisotropic group-scheme generalization and an existence statement for orders. Several hypotheses and the targets of these extensions require clarification. This is a statement-only draft requiring human mathematical review.
\subsection{Sources}
\begin{itemize}
\item Local UnsolvedMath registry, permanent ID \(30003793\), title ``Descent of Hereditary Orders with Involution,'' exact source problem code \texttt{catalogue entry}; supplied collection: Oberwolfach Reports Open Problems. No individual proposer is named. The supplied proposed year is 2018.
\item Supplied report citation: \emph{Quadratic Forms and Related Structures over Fields (2019)}, Oberwolfach Reports, DOI \texttt{10.4171/owr/2018/21}, pp.\ 1235--1291. \url{https://doi.org/10.4171/owr/2018/21}. The supplied title year is 2019, while the DOI contains 2018; this discrepancy has not been resolved.
\item Supplied literature source, labeled ``orders with involution'': \url{https://arxiv.org/abs/1804.09527}.
\item Supplied literature source, labeled ``valuations on algebras with involution'': \url{https://doi.org/10.1007/s00208-010-0580-9}.
\item Supplied literature source, labeled ``workshop report'': \url{https://ems.press/content/serial-article-files/46742}.
\item Literature assessment dated exactly \texttt{2026-08-21}: the supplied curation labels the problem open and describes the anisotropic or residually-anisotropic strengthening as a conjecture. It states that related valuation and Grothendieck--Serre results do not prove this exact hereditary-order assertion. This is supplied dated curation derived from the OWR record, not an independently verified literature assessment; the status and cited claims require independent verification.
\item The registry labels the statement \texttt{corrected\_verified} and reports a manual semantic, typographic, and LaTeX review. These are supplied review claims and do not establish mathematical correctness or resolve the qualifications identified here.
\end{itemize}
\end{document}
